\needspace{12\baselineskip}
\begin{center}
\pdfbookmark[1]{8. On Art. 114 and 177.}{zusatz-008}
[8. On Art. 114 and 177.]
\end{center}
\begin{center}
\{Misprints in Dr. \textsc{Gauss}' Theoria motus corporum coelestium etc. Hamburgi 1809, by Senator Bar. \textsc{Oriani}.\}
\end{center}
\begin{center}
\begin{tabular}{c}\hline
\small Monatliche Correspondenz Volume XXI, pp. 282--283, March 1810.\\\hline
\end{tabular}
\end{center}
\{.... For equation 7 [art. 114] to be an identity, the coefficients of \(\delta\delta'\delta''\), \(D\delta'\delta''\), \(DD'\delta''\) etc. must be \(=0\); but after the appropriate reductions the coefficient of \(\delta\delta'\delta''\) is
\SourceLeaf{308}
\[
(0.1.2)=\left\{\begin{aligned}
&\operatorname{tang}\beta'\operatorname{tang}\beta\,\sin(L''-\alpha'')\sin(L'-L)\\
&+\operatorname{tang}\beta'\operatorname{tang}B\,\sin(L''-\alpha'')\sin(L'-\alpha)\\
&+\operatorname{tang}\beta\operatorname{tang}\beta''\sin(L'-\alpha')\sin(L-L'')\\
&+\operatorname{tang}B\operatorname{tang}\beta''\sin(L'-\alpha')\sin(\alpha-L'')\\
&+\operatorname{tang}\beta''\operatorname{tang}\beta'\sin(L-\alpha)\sin(L''-L')\\
&+\operatorname{tang}\beta''\operatorname{tang}B'\sin(L-\alpha)\sin(L''-\alpha')\\
&+\operatorname{tang}B'\operatorname{tang}\beta\,\sin(L''-\alpha'')\sin(\alpha'-L)\\
&+\operatorname{tang}B'\operatorname{tang}B\,\sin(L''-\alpha'')\sin(\alpha'-\alpha)\\
&+\operatorname{tang}\beta\operatorname{tang}B''\sin(L'-\alpha')\sin(L-\alpha'')\\
&+\operatorname{tang}B\operatorname{tang}B''\sin(L'-\alpha')\sin(\alpha-\alpha'')\\
&+\operatorname{tang}B''\operatorname{tang}\beta'\sin(L-\alpha)\sin(\alpha''-L')\\
&+\operatorname{tang}B''\operatorname{tang}B'\sin(L-\alpha)\sin(\alpha''-\alpha').
\end{aligned}\right.
\]
Hence, if we have only
\[
\left.\begin{aligned}
&\operatorname{tang}\beta'\operatorname{tang}\beta\,\sin(L''-\alpha'')\sin(L'-L)\\
&+\operatorname{tang}\beta\operatorname{tang}\beta''\sin(L'-\alpha')\sin(L-L'')\\
&+\operatorname{tang}\beta''\operatorname{tang}\beta'\sin(L-\alpha)\sin(L''-L')
\end{aligned}\right\}=0,
\]
the same coefficient does not become \(=0\), but it is further required that
\[
B=0,\qquad B'=0,\qquad B''=0.\}
\]
\begin{center}\rule{.25\textwidth}{.35pt}\end{center}

\{The elegant theorem attributed [art. 117] to the illustrious \textsc{Laplace} was in fact first discovered by \textsc{Leonardo Eulero}. For in Comment. Acad. Petropol. Tom. XVI\HistoricalNoteEight{Manuscript remark by \textsc{Gauss}:] \textls[150]{This} theorem is not found in the cited work, but the following is found on p. 101: \(\int_0^1\BareRoot{\left(\log\frac1x\right)}\mathrm{d}x=\tfrac12\BareRoot{\pi}.\) If we put \(x=e^{-tt}\) here, we obtain \(\int_0^\infty2tte^{-tt}\mathrm{d}t=\tfrac12\BareRoot{\pi}.\) But \(2tte^{-tt}\mathrm{d}t=-\mathrm{d}(te^{-tt})+e^{-tt}\mathrm{d}t\) and \(te^{-tt}=0\) both for \(t=0\) and for \(t=\infty\), so \(\int_0^\infty e^{-tt}\mathrm{d}t=\tfrac12\BareRoot{\pi}.\)} \textsc{Eulerus} shows that the integral \(-\int\frac{\mathrm{d}x}{\BareRoot{\left(\log\frac1x\right)}}\), taken from \(x=1\) to \(x=0\), is \(=\BareRoot{\pi}\), where \(\pi\) is the semicircumference of a circle drawn with radius \(=1\). Now, on putting \(x=e^{-tt}\), we have
\[
\frac{-\mathrm{d}x}{\BareRoot{\left(\log\frac1x\right)}}=2e^{-tt}\mathrm{d}t.
\]

Therefore the integral \(\int e^{-tt}\mathrm{d}t\) from \(t=0\) to \(t=\infty\) will be \(=\tfrac12\BareRoot{\pi}\), and consequently the same integral from \(t=-\infty\) to \(t=+\infty\) will become \(=\BareRoot{\pi}.\)\}
\begin{center}\rule{.25\textwidth}{.35pt}\end{center}

\SourceLeaf{309}
\begin{center}
\begin{tabular}{c}\hline
\small Monatliche Correspondenz. Volume XXI. P. 280. March 1810.\\\hline
\end{tabular}
\end{center}
I am much obliged to Mr. \textsc{Oriani} for noting the misprints in my Theoria. He is quite right that on p. 129 [Art. 114] I forgot to add that \(B,B',B''=0\) must be assumed if the condition equation under which equation 7 is unusable is to have the form given there. It is, moreover, clear that even if \(B,B',B''=0\) does not hold, equation 7 can still be unusable, namely if the expression of 12 terms developed by \textsc{Oriani} happens to become \(=0\) or very small.

It had already occurred to me earlier that \textsc{Euler} had found the theorem from which the beautiful theorem I attributed to \textsc{Laplace} can very easily be derived, but the passage on p. 212 [Art. 177] had already been printed; I did not wish to put this among the errata, however, because \textsc{Laplace} was at least the first to state the above theorem in the form used there. ....
\begin{center}\rule{.18\textwidth}{.35pt}\end{center}
